Gazeboko \texttt{simple} ingurunean probak etengabeak izan ziren garapenean zehar, parametro bat aldatzen genuen bakoitzean (atalaseak, abiadurak, noraeza...) robotaren portaera berriz egiaztatu behar baitzen. Iraupen desberdinekin probatu genuen (5, 7 eta 10 minutu): garapeneko iterazio azkarretarako 5 minutuko probak erabili genituen doikuntzak arinago balioztatzeko, eta behin betiko portaera ebaluatzeko 10 minutuko saio osoak egin genituen hainbat hasierako posiziotatik.

Hasierako posizioen proba osagarrietan, erdigunetik abiatzean 17 segundo inguru behar izan zituen pasabidean sartzeko; pasabidearen aurre-aurretik hasita, ordea, ez zen sartu, biratzean bideraketa desbideratu zelako. Horrek erakusten du aurreko ibilbideak eta noraezaren dinamikak hasierako hurbiltasun fisikoak baino pisu handiagoa dutela gelara sartzeko unean.
